\section{Quasilocal Algebra}\label{sctn:quasilocal-algebra}
The second point that requires clarification is a bit more subtle.

When one requires that $\mathfrak{U}(\mathbf{B_1})$ and $\mathfrak{U}(\mathbf{B_2})$ commute, this tacitly implies that $\mathfrak{U}(\mathbf{B_1})$ and $\mathfrak{U}(\mathbf{B_2})$ are subalgebras of a single, larger algebra, an algebra in which products and sums of their elements are defined.

Isotony is the natural tool for this. If, say, $\mathbf{B_1} \subset \mathbf{B_2}$, isotony
places $\mathfrak{U}(\mathbf{B_1})$ inside $\mathfrak{U}(\mathbf{B_2})$, so that both live in the
single algebra $\mathfrak{U}(\mathbf{B_2})$. That case is of no use here, since nonempty nested regions
are never completely spacelike. But isotony does not require one of the two regions to contain the
other: it suffices to find a \emph{third} basis element $\mathbf{B}$ containing both. In Minkowski
spacetime such a $\mathbf{B}$ always exists --- any two double cones $I^+(p_1) \cap I^-(q_1)$ and
$I^+(p_2) \cap I^-(q_2)$ lie inside a larger double cone $I^+(p) \cap I^-(q)$, by taking $p$
far enough in the past of $p_1$ and $p_2$ and $q$ far enough in the future of $q_1$ and $q_2$.
Isotony then places $\mathfrak{U}(\mathbf{B_1})$ and $\mathfrak{U}(\mathbf{B_2})$ inside the single,
larger algebra $\mathfrak{U}(\mathbf{B})$, in which products and sums of their elements are
defined. We can thus interpret the statement

\begin{quote}
    If $\mathbf{B_1}$ and $\mathbf{B_2}$ are completely spacelike with respect to each other, then $\mathfrak{U}(\mathbf{B_1})$ and $\mathfrak{U}(\mathbf{B_2})$ commute.
\end{quote}

as the statement

\begin{quote}
    If $\mathbf{B_1}$ and $\mathbf{B_2}$ are completely spacelike with respect to each other, then $\mathfrak{U}(\mathbf{B_1})$ and $\mathfrak{U}(\mathbf{B_2})$ commute in $\mathfrak{U}(\mathbf{B})$ for every basis element $\mathbf{B}$ containing both $\mathbf{B_1}$ and $\mathbf{B_2}$.
\end{quote}

Two remarks. First, the choice of $\mathbf{B}$ does not matter: because the isotony embeddings
are injective and compatible, commutation in one containing algebra implies commutation in every
other (using a basis element containing both of them, which exists as above). Second, the quasilocal algebra $\mathfrak{U}$ of the next chapter contains every
$\mathfrak{U}(\mathbf{B})$, so the statement can equally be read as commutation in
$\mathfrak{U}$; the two readings are equivalent. The formulation through a containing
$\mathfrak{U}(\mathbf{B})$ is preferred because it uses only the net and isotony, and because it
is the formulation that survives in curved spacetime (Chapter 9), where a containing basis element
need not exist and the quasilocal algebra need not exist either.

Clarification.
